\chapter{Reproducibility Guide}\label{app:reproducibility}

\section{Where everything is}
\begin{itemize}
\item \textbf{Code, configurations and raw logs:} the project repository \url{https://github.com/Unavailable4u/meeting-loss} (commit \code{\repocommit} for the analysis in this report). Raw run logs are in \code{results/raw/} (one JSON-lines file per run, with a header that records the configuration, its hash, the git commit and whether the working tree was clean); processed tables are in \code{results/processed/}.
\item \textbf{Locks:} tags \code{prereg-v1} (\code{3378d36}), \code{prereg-v2} (\code{cc2830c}) and \code{prereg-v3} (\code{a49e9e3}). No source, configuration or hypothesis file locked by a tag changed between the tag and the final commit, except where the drift log records an additive change.
\item \textbf{This report:} the \code{scripts/} folder contains the Python code that produces every figure (\code{fig\_study1.py}, \code{fig\_study23.py}) and every table (\code{make\_tables.py}) directly from the raw logs, using the repository's own scoring and bootstrap functions (\code{src/sweep.py}, \code{src/prereg3.py}). The \LaTeX{} tables in \code{tables/} and the PDF figures in \code{figures/} are the outputs of these scripts, so the numbers in the text can be re-checked.
\end{itemize}

\section{Regenerating figures and tables}
\begin{lstlisting}
git clone https://github.com/Unavailable4u/meeting-loss.git
cd meeting-loss
pip install numpy pandas pyyaml matplotlib
cd report/scripts
python make_tables.py       # writes ../tables/*.tex
python fig_study1.py        # writes ../figures/fig_s1_*.pdf, fig_tuning_noise, fig_dose, ...
python fig_study23.py       # writes ../figures/fig_s2_*.pdf, fig_s3_*.pdf, fig_metrics.pdf
cd .. && latexmk -pdf main.tex
\end{lstlisting}
The scripts find the repository root automatically; the environment variable \code{MEETING\_LOSS\_REPO} overrides it.
The figure scripts import the repository's \code{src} package; \code{src/prereg3.py} needs only NumPy and the standard library at analysis time. Re-running the experiments themselves requires PyTorch (CPU is sufficient) and the commands in \code{scripts/README.md} of the repository; the 260 confirmatory runs of Study~3 took 15\,856 seconds ($4.4$ hours) on four worker processes, as logged.

\section{Seed ranges}
\begin{table}[h]
\centering\small
\caption{Seed ranges by purpose. No seed is used for both tuning and evaluation, and the exploratory seeds are never reused in a confirmatory test.}
\begin{tabular}{@{}l l@{}}
\toprule
Purpose & Seeds\\
\midrule
Tuning (selection), tags \code{t1}, \code{t2}, \code{t5}--\code{t8} & 1000/1001, 2000/2001, 5000/5001, \dots, 8000/8001\\
Tuning of B48, S30T, SBT (\code{t9}, \code{t10}, \code{t11}) & 9000/9001, 9100/9101, 9200/9201\\
Study 1 (confirmatory) & 0--19 (A, D, B, A2); 0--9 (other groups)\\
Concentration, geometry, pilot (exploratory) & 500--502; 600--602\\
Dose-response blocks & 1100--1109, 1110--1119, 1120--1129\\
Timing and small-batch looks & 1200--1204\\
Study 2 (confirmatory) & 1300--1319\\
Headroom and look 2 (exploratory) & 2100--2109, 2110--2119\\
Study 3 (confirmatory) & 2300--2319\\
\bottomrule
\end{tabular}
\end{table}

\section{Status of verification at the time of writing}
\begin{itemize}
\item \textbf{Numbers.} Every number in the report tables and figures was recomputed from the raw logs by the scripts above; the recomputed values reproduce the values in the repository's experiment log for Study 1 (E-008), Study 2 (E-022) and Study 3. The repository's own documents record the verdicts of H8--H10 (\code{docs/HYPOTHESES.md}, E-027 in \code{docs/EXPERIMENT\_LOG.md}, F-025 in \code{docs/FINDINGS.md}, drift entry DR-018); they were added in commit \code{c6d8ce5}, after the analysis commit, and the recomputation here gives the same verdicts.
\item \textbf{Unit tests.} The repository contains 24 test files with 223 test functions, which pytest expands to 244 test cases. All 244 passed on 25 November 2026 (Python 3.13.18, PyTorch 2.15.0, NumPy 2.5.5, CPU only, 36~s): \code{pytest tests/}.
\item \textbf{Bibliography.} 29 of the 31 entries were checked against web sources (publisher, proceedings, arXiv or DOI registry pages) on 25 November 2026: authors and their order, title, venue, year, and volume and page range where a page gave them; the other two (Lee et al., DTP, and Kaddour et al.) were checked on 31 October 2026. Corrections made in that check are expanded author lists (the PyTorch, Benchmarking Neural Network Training Algorithms, Accounting for Variance and Prioritized Training papers), the ICLR author order of \citet{dehghani2022efficiency}, and added volumes and page ranges. Page ranges that no fetched page gave are left out (NeurIPS papers, the PyTorch paper). The literature review is a focused one and the systematic related-work check listed in \code{docs/LITERATURE.md} remains open (\cref{sec:future}).
\end{itemize}
